\chapter{补充推导与公式核验}\label{app:proofs}
本附录补全正文没有展开到底的步骤，也提供几种独立的符号检查方式。它与第二章使用同一时间约定。

\section{ELBO 的望远镜重排}
直接由负下界得到
\begin{align}
 L=&\E_q\left[-\log p(x_K)-\sum_{k=1}^K\log p_\theta(x_{k-1}\mid x_k)
 +\sum_{k=1}^K\log q(x_k\mid x_{k-1})\right].
\end{align}
对 $k\ge2$，贝叶斯公式给出
\begin{equation}
 \log q(x_k\mid x_{k-1})=
 \log q(x_{k-1}\mid x_k,x_0)+\log q(x_k\mid x_0)-\log q(x_{k-1}\mid x_0).
\end{equation}
最后两个对数对 $k$ 求和会相互抵消，留下 $\log q(x_K\mid x_0)-\log q(x_1\mid x_0)$。再与 $k=1$ 的 $\log q(x_1\mid x_0)$ 抵消，便得到正文式\eqref{eq:elbo}的三个部分。每个 KL 的条件变量和期望分布必须一并保留；只写一个抽象的“KL 求和”容易漏掉重建项。

\section{经典 score matching 的积分分部}
在密度 $p$ 足够光滑、边界项消失的情况下，考虑
\begin{equation}
 J(\theta)=\tfrac12\E_p\norm{s_\theta(X)-\grad\log p(X)}^2.
\end{equation}
展开并略去与参数无关的部分，
\begin{equation}
 J=\tfrac12\int p\norm{s_\theta}^2\dd x
 -\int s_\theta\cdot\grad p\dd x+C.
\end{equation}
逐坐标积分分部得到
\begin{equation}
 -\int s_\theta\cdot\grad p\dd x=\int p\divg s_\theta\dd x.
\end{equation}
所以
\begin{equation}
 J(\theta)=\E_p\left[\tfrac12\norm{s_\theta(X)}^2+\divg s_\theta(X)\right]+C.
\end{equation}
未知的真实 score 从可优化项中消失。这是经典 score matching 的关键结构\cite{r_scorematching}。但高维网络需要计算散度，未加噪数据也可能集中在低维流形。去噪 score matching 改用已知平滑核的梯度，不需要直接计算这项散度，并在多个噪声尺度学习平滑后的场。

\section{反向漂移的局部贝叶斯计算}
先假设 $g$ 仅依赖时间且局部看作常数，$f$ 和密度足够光滑。令 $h>0$，给定当前状态 $X_t=x$，前一状态写作 $y=x+\delta$。小步转移密度近似为
\begin{equation}
 q_h(x\mid y)\propto\exp\left[-\frac{\norm{x-y-hf(y,t-h)}^2}{2hg^2}\right].
\end{equation}
逆条件密度正比于 $p_{t-h}(x+\delta)q_h(x\mid x+\delta)$。主导的 $\delta$ 尺度为 $\sqrt h$。对对数先验作展开：
\begin{equation}
 \log p_{t-h}(x+\delta)
 =\log p_t(x)+s_t(x)\cdot\delta+O(h)+O(\norm\delta^2).
\end{equation}
保留决定均值一阶项的部分，转移指数为
\begin{equation}
 -\frac{\norm{\delta+hf(x,t)}^2}{2hg^2}
 +s_t(x)\cdot\delta.
\end{equation}
配方后，均值为 $-hf(x,t)+hg^2s_t(x)$，协方差主导项为 $hg^2I$。漂移导数和密度二阶项对均值、协方差给出更高阶修正；使该展开严格需控制余项和尾部。该计算展示了反向 drift 的来源，不代替一般反向扩散定理的完整条件证明。

\section{矩阵扩散且依赖状态时的修正}
对于
\begin{equation}
 \dd X_t=f(X_t,t)\dd t+G(X_t,t)\dd W_t,
 \qquad A(x,t)=G(x,t)G(x,t)^\top,
\end{equation}
Fokker--Planck 方程为
\begin{equation}
 \partial_tp=-\sum_i\partial_i(f_ip)
 +\tfrac12\sum_{i,j}\partial_i\partial_j(A_{ij}p).
\end{equation}
定义向量 $(\divg A)_i=\sum_j\partial_jA_{ij}$。将二阶项展开为概率流，得到概率流速度
\begin{equation}
 v^{\mathrm{PF}}=f-\tfrac12\divg A-\tfrac12A\grad\log p.
\end{equation}
使用 $t$ 向后走的反向 SDE 漂移则为
\begin{equation}
 f^{\mathrm{reverse}}=f-\divg A-A\grad\log p.
\end{equation}
相应的逆向正时钟漂移要整体取负号。正文 $A=g(t)^2I$ 不依赖状态，$\divg A=0$，才化成熟悉的标量公式。把状态依赖扩散硬套进 $f-g^2s$ 会漏掉散度修正。这一点可对照 Score SDE 原文附录 A\cite{r_sde}。

\section{通过流映射的 Jacobian 检查密度符号}
设 $T_t(z)$ 是 ODE 流，$J_t=\partial_zT_t$。对初始状态求导：
\begin{equation}
 \dot J_t=(\partial_xv_t)(T_t)J_t.
\end{equation}
利用行列式的导数公式，
\begin{equation}
 \frac{\dd}{\dd t}\log|\det J_t|
 =\tr(J_t^{-1}\dot J_t)=\divg v_t(T_t).
\end{equation}
变化量公式 $p_t(T_t(z))|\det J_t(z)|=p_0(z)$ 因而给出
\begin{equation}
 \frac{\dd}{\dd t}\log p_t(T_t(z))=-\divg v_t(T_t(z)).
\end{equation}
空间体积膨胀时密度降低，这与正文的连续性方程推导一致。这个证明要求可微可逆流；奇异端点则需要单独处理，不能直接使用可逆变化量公式。

\section{一维高斯概率流的完整解析轨迹}
设边缘路径为 $p_t=\N(m_t,h_t)$，$h_t>0$。选择速度
\begin{equation}
 v_t(x)=m_t'+\frac{h_t'}{2h_t}(x-m_t).
\end{equation}
令 $x_t=m_t+\sqrt{h_t}\,z$，直接求导即验证它满足此 ODE。若 $z\sim\N(0,1)$，边缘恰为 $p_t$。因此给定任意可微高斯均值和正方差路径，生成速度由“移动中心”与“缩放偏离中心的位置”两项组成。

对于 VP 前向过程，$m_t=a(t)m_0$、$h_t=a(t)^2h_0+b(t)^2$。取 $s_t=-(x-m_t)/h_t$，代入 $f-g^2s/2$，利用 $m_t'=\gamma m_t$ 与 $h_t'=2\gamma h_t+g^2$，得到同一速度。这个独立检查能同时核验 score 符号和概率流的半系数。

\section{CFM 等梯度结论的精确范围}
记 $V$ 为条件速度标签，$U=(Y_\tau,\tau)$，$m(U)=\E[V\mid U]$。由投影恒等式，
\begin{equation}
 L_{\mathrm{CFM}}(\theta)-L_{\mathrm{FM}}(\theta)
 =\E\norm{V-m(U)}^2.
\end{equation}
如果配对分布与概率路径固定，不随 $\theta$ 改变，右边为常数；在可交换求导与积分的条件下，两者梯度相同。若用当前模型生成配对，再让配对的采样过程也参与反向传播，右边一般会随参数变化。reflow 等过程常把新配对视为固定训练数据，再训练下一轮模型；不能忽略“路径是否固定”这个条件。

类似地，可更换时间抽样或正权重，但必须保证比较的两个损失使用相同权重。模型容量不足时，“每个位置的理论最优相同”也不意味着两种有限网络训练具有完全相同的优化轨迹。

\section{热方程与熵：加噪的一个背景直觉}
纯布朗加噪 $X_t=X_0+W_t$ 的密度满足 $\partial_tp_t=\Delta p_t/2$。在密度、尾部和导数足以支持以下运算时，微分熵 $H(p_t)=-\int p_t\log p_t$ 的变化为
\begin{align}
 \frac{\dd H}{\dd t}
 &=-\tfrac12\int(1+\log p_t)\Delta p_t\dd x\\
 &=\tfrac12\int\frac{\norm{\grad p_t}^2}{p_t}\dd x
 =\tfrac12\E_{p_t}\norm{s_t(X_t)}^2\ge0.
\end{align}
右边称为与平移有关的 Fisher 信息的一种形式。纯加性噪声使分布变得更不集中，而 score 的平方平均描述此熵增长速率。对于带漂移的 VP 过程，熵导数还有漂移散度贡献；不能将这条单调性无条件移植到所有扩散过程。

